The number of rotary channels does not uniquely determine the positional
responses available to an attention head. We study whether frequency allocation
can reduce the cost of a smaller rotary budget while keeping total query/key
width, model size, and training resources fixed. A finite-distance analysis bounds
the spectral tail contributed by nearly repeated slow-frequency subspaces.
Under a bounded coefficient norm, this limits responses outside their shared
span; a controlled softmax result states the additional conditions needed to
turn that restriction into distributional error. These are conditional basis
results, not language-model loss bounds. We instantiate the question with a
fixed-support anchored EVQ-Cosh allocation and a six-arm factorial design:
32 versus 16 rotary pairs, geometric versus analytic allocation, and two
16-pair frequency-subset controls. Existing fixed-support training and single-budget
MLA results motivate the experiment but do not identify a budget interaction.
The new paired training study is in progress. Whether analytic allocation
preserves long-context quality at half the rotary width, without impairing the
training window or losing to simple subset selection, remains an empirical
question.
